\subsection{Molecular Replacement for Cryo-EM Maps}

A molecular replacement (MR) pipeline has been implemented for
local fitting known structures into cryo-EM density maps.
We have previously discussed the implementation of "Jiggle Fit"
\cite{casanal-protein-science} to address the same problem.
Here, the method is extended with the hope of increasing
the robustness of the method by increasing the distance
between initial centre and true position.

Discuss Molrep, EM Placement, Situs, see Google Docs EM notes

This method adapts
the fast rotation function of Crowther~\cite{Crowther1972} to operate
on real-space density maps rather than crystallographic structure
factors, combined with a phased translation function \cite({} and
density-weighted recentring to handle the difference between user input
and user input. The implementation is designed for interactive use:
the user identifies a region of unmodelled density on screen and supplies
a search model, and the pipeline creates a new molecule for the
top scoring solution.


\subsubsection{Density-Weighted Recentring}

In practice the user's click will be offset from the true centre of
mass of the target molecule. Because the Crowther rotation function
employs a Tukey (tapered cosine) window that must be symmetric about
the molecule centre, an asymmetric offset corrupts the Patterson map
and degrades the rotation signal. To correct for this, the pipeline
performs iterative density-weighted recentring before extracting the
search fragment.

A spherical region of radius $r_c = 0.75\,R_g$ (where $R_g$ is the
radius of gyration of the search model) is extracted around the
current centre estimate. The density statistics (mean $\bar{\rho}$ and
standard deviation $\sigma$) of the extracted region are computed, and
the density-weighted centroid of voxels exceeding a threshold of
$\bar{\rho} + 3\sigma$ is calculated:
%
\begin{equation}
  \Delta\mathbf{r} = \frac{\sum_{i} \rho_i\,\mathbf{r}_i}
                          {\sum_{i} \rho_i}
  \quad\text{where } \rho_i > \bar{\rho} + 3\sigma
\end{equation}
%
The shift per iteration is clamped to 3~\AA\ to prevent oscillation,
and the total accumulated shift is limited to $R_g$ to prevent drift
toward neighbouring molecules. Iteration terminates when the shift
falls below 0.1~\AA. The small centroid radius relative to the
fragment extraction radius ensures that only density from the target
molecule contributes to the centroid, avoiding contamination from
neighbouring molecules in the map.

\subsubsection{Crowther Fast Rotation Function}

The rotation function follows the formulation of
Crowther~\cite{Crowther1972} as extended by Navaza~\cite{Navaza1994}.
The target and model Patterson maps are expanded in spherical harmonics
$Y_l^m(\theta, \phi)$ on concentric radial shells:
%
\begin{equation}
  a_{lm}(r) = \sum_{\mathbf{s}} |F(\mathbf{s})|^2\, j_l(2\pi s r)\,
    Y_l^{m*}(\hat{\mathbf{s}})
\end{equation}
%
where $j_l$ is the spherical Bessel function of the first kind and the
sum runs over the structure factor reflections obtained by Fourier
transformation of the extracted density fragments.

The cross-correlation coefficients are computed by radial integration:
%
\begin{equation}
  c_{l,m,m'} = \int_0^R a_{lm}^*(r)\, b_{lm'}(r)\, r^2\, \mathrm{d}r
\end{equation}
%
where the integration is discretised over 50 radial shells up to the
Crowther radius $R = 0.9\,R_\text{frag}$, with
$R_\text{frag} = 1.5\,R_g$.

The rotation function is then evaluated as a function of ZYZ Euler
angles $(\alpha, \beta, \gamma)$:
%
\begin{equation}
  \mathcal{R}(\alpha, \beta, \gamma) = \sum_{l=2}^{l_\text{max}}
    \sum_{m=-l}^{l} \sum_{m'=-l}^{l}
    c_{l,m,m'}\, d^l_{mm'}(\beta)\,
    e^{-im\alpha}\, e^{-im'\gamma}
\end{equation}
%
where $d^l_{mm'}(\beta)$ are the Wigner $d$-matrix elements. The
double exponential sum in $\alpha$ and $\gamma$ is evaluated
efficiently as a two-dimensional discrete Fourier transform for each
sampled value of $\beta$, making the computational cost
$\mathcal{O}(N_\beta \cdot N^2 \log N)$ where $N = 2l_\text{max}+1$.
The implementation uses $l_\text{max}=15$ and $N_\beta=36$.

Both the target and model density fragments are extracted with a Tukey
window of parameter $\alpha_T = 0.5$, providing a flat response from
the centre to $0.5\,R_\text{frag}$ with a cosine taper to zero at
$R_\text{frag}$, ensuring that both Pattersons are comparably
windowed.

\subsubsection{Rotation Refinement and Filtering}

The coarse rotation search produces peaks on a grid with
$\sim\!5^\circ$ spacing. Peaks scoring below 70\% of the top score are
discarded. The surviving peaks undergo two rounds of local refinement
on a $3\times3\times3$ grid of ZYZ Euler angles, first with a
$2^\circ$ step and then with a $0.7^\circ$ step. After refinement,
solutions scoring below 80\% of the top refined score are removed.
Finally, a proximity filter removes solutions within $4^\circ$ of a
higher-scoring solution (measured as the geodesic angle between unit
quaternions, $\theta = 2\cos^{-1}|\mathbf{q}_1 \cdot \mathbf{q}_2|$),
eliminating near-duplicate orientations that arise from the finite
grid sampling.

\subsubsection{Phased Translation Function}

For each surviving rotation, the search model is rotated and its
electron density is computed on the same grid as the observed map
using Clipper's \texttt{EDcalc\_iso} calculator. A phased translation
function is then evaluated by cross-correlating the observed and
calculated maps in reciprocal space:
%
\begin{equation}
  T(\mathbf{t}) = \mathcal{F}^{-1}\!\left[
    F_\text{obs}(\mathbf{h})\, F_\text{calc}^*(\mathbf{h})
  \right]
\end{equation}
%
The resulting real-space map is scanned for peaks, which are scored in
units of standard deviations ($\sigma$) above the mean. Peaks are
filtered to retain only those within $R_\text{frag}$ of the
(recentred) target position, and solutions with translation scores
below 70\% of the top score are discarded.

\subsubsection{Solution Scoring}

Each surviving solution is scored by computing the mean observed
density at the C$\alpha$ positions (or C1$'$ for nucleic acids) of
the placed model, sampled by cubic interpolation from the original
(unmasked) cryo-EM map. This provides an independent validation of
solution quality that does not depend on the rotation or translation
function scores.

\subsubsection{Parallelisation}

The pipeline exploits thread-level parallelism at three stages: the
radial shell integration during coefficient computation, the local
rotation refinement of individual peaks, and the per-rotation
translation searches (each comprising map calculation, FFT correlation
and peak filtering). The translation search stage, which dominates the
overall computation time, achieves near-linear speedup with the number
of available cores.

\begin{thebibliography}{9}

\bibitem{Crowther1972}
Crowther, R.\,A. (1972).
\textit{The Molecular Replacement Method}
(ed.\ M.\,G.\ Rossmann), pp.\ 173--178. New York: Gordon \& Breach.

\bibitem{Navaza1994}
Navaza, J. (1994).
\textit{AMoRe}: an automated package for molecular replacement.
\textit{Acta Cryst.}\ \textbf{A50}, 157--163.

\end{thebibliography}
